\section{Experimental Studies}
\label{sec:exp}

The experiments evaluated overall optimisation performance, the contribution
of the two modules, and the group and batch properties associated with their
design. The main comparisons used 10 and 20 objectives. Additional objective
counts and separate mechanism studies provided complementary evidence.

\subsection{Experimental Settings}
\label{sec:exp:settings}

\paragraph{Benchmarks and evaluation budgets.}
The main test suite comprised DTLZ1--DTLZ7 and WFG1--WFG9. Both the 10- and
20-objective comparisons used a population setting of $N=100$ and a budget of
$FE_{\max}=300$ real evaluations. The recorded decision dimension was $D=30$,
except for WFG3, which used $D=31$. Additional results covered $M=3,5,8$ on the
same 16 problems and $M=12$ on WFG1--WFG9. Table~\ref{tab:exp:settings}
distinguishes these settings from the mechanism experiments. In particular,
the grouping studies used 500 evaluations, whereas the candidate-value probe
used 300. Their results were analysed within their respective protocols.

\begin{table*}[t]
\centering
\caption{Experimental settings of the retained result sets. A dash denotes unverified or unavailable metadata.}
\label{tab:exp:settings}
\footnotesize
\setlength{\tabcolsep}{3pt}
\renewcommand{\arraystretch}{1.12}
\begin{tabular}{lccccc}
\toprule
Result set & $M$ & $D$ & $N$ & $FE_{\max}$ & Runs/config. \\
\midrule
Main comparison & 10, 20 & 30 (31) & 100 & 300 & -- \\
Additional full suite & 3, 5 & 10 & -- & -- & -- \\
Additional full suite & 8 & 10 (11) & 100 & 300 & -- \\
WFG-only comparison & 12 & 30 (31) & 100 & 300 & -- \\
Two-module ablation & 10, 20 & 30 (31) & 100 & 300 & -- \\
Snapshot/group audit & 10, 20 & 30 (31) & 100 & 500 & 5/behaviour \\
Good-group precision & 10, 20 & 30 (31) & 100 & 500 & 25 \\
Candidate-value probe & 20 & 30 (31) & 100 & 300 & 10/policy \\
\bottomrule
\end{tabular}
\par\smallskip
\begin{minipage}{\linewidth}
\scriptsize Parenthesised decision dimensions apply to WFG3. The snapshot audit has two behaviours per configuration; the candidate probe has five policies. The good-group precision results cover the original five-problem study. Spreadsheet run counts are intentionally left unfilled.
\end{minipage}
\end{table*}

\paragraph{Algorithms and indicators.}
HPDC-MaOEA was compared with REMO, PIEA, MCEAD, R2AEA, KRVEA and CSEA at both
$M=10$ and $M=20$, with PCSAEA additionally available at $M=10$ and HES-EA at
$M=20$. REMO, PIEA, MCEAD and KRVEA formed the common subset across the
additional objective counts. All available baselines were retained in each
main comparison. The implementation recorded as
\texttt{REMO\_new2\_AdaMaO\_SDEOnly\_UniformMix\_Original} is denoted
HPDC-MaOEA in the performance tables.
The primary indicator was inverted generational distance (IGD), for which
smaller values indicate better approximation quality. Entries report the
recorded mean and standard deviation; the lowest mean in each row is bold.
The separate routing study also supplied IGD$^+$ results at $M=20$.

\paragraph{Statistical reporting.}
The main comparison exports identify their pairwise annotations as Wilcoxon
rank-sum results. We preserve their direction: $+$ denotes a comparator
reported as better than HPDC-MaOEA, $-$ a comparator reported as worse, and
$=$ no detected difference. Summary counts are expressed from HPDC-MaOEA's
perspective as wins/ties/losses (W/T/L), so an exported $-$ contributes a win.
Counts were reconstructed from the individual cells. These annotations are
reported as supplied; no new significance test was inferred from means and
standard deviations. In the mechanism studies, repeated checkpoints were
first aggregated within each run. The grouping complementarity tests used
25 runs per problem--objective--stage cell and Holm adjustment over 40 cells.
The candidate probe is reported through run-level descriptive statistics.

\TODO{Complete the independent-run counts, original test settings and
multiple-comparison policy for the spreadsheet experiments; recover the
missing $N$ and FE settings for $M=3,5$, and the execution platform and IGD
reference-set specifications. Verify the identical PCSAEA/KRVEA entries for
DTLZ1 at $M=10$ before finalising the corresponding comparisons.}

\subsection{Comparison with State-of-the-Art Algorithms}
\label{sec:exp:comparison}

HPDC-MaOEA obtained more reported wins than losses across the available
comparators in both main settings. The aggregate W/T/L counts were
73/18/21 at $M=10$ and 70/21/21 at $M=20$, each over 112 problem--comparator
pairs (Tables~\ref{tab:exp:main10} and~\ref{tab:exp:main20}). Against REMO,
the counts were 9/5/2 and 8/6/2, respectively; against CSEA, they were
12/3/1 and 11/4/1. The comparison with PIEA was more problem dependent,
with 6/3/7 at $M=10$ and 9/2/5 at $M=20$.

The strongest common improvements over REMO occurred on DTLZ7 and several
WFG problems. On DTLZ7, mean IGD decreased from 21.252 to 12.122 at $M=10$
and from 38.239 to 18.154 at $M=20$, reductions of 43.0\% and 52.5\%.
HPDC-MaOEA also attained the lowest mean among the available algorithms on
WFG7 in both settings, and on WFG1 and WFG2 at $M=10$. These results show
competitive performance across several many-objective instances under a
300-evaluation budget.

The problemwise results also distinguish this advantage from uniform
dominance. All seven comparators had lower reported IGD on WFG3 in both
main settings. R2AEA achieved lower mean IGD on DTLZ7 at both objective
counts, while PIEA and MCEAD were stronger on DTLZ5 and DTLZ6. Across the
four common baselines, the median relative IGD reduction was 10.49\% at
$M=10$ and 9.73\% at $M=20$; the corresponding arithmetic means were
3.37\% and $-0.17\%$. Relative reduction was calculated per
problem--comparator pair as $(\mathrm{IGD}_{b}-\mathrm{IGD}_{H})/
\mathrm{IGD}_{b}$, using the exported means. The difference between the
median and mean reflects the uneven magnitude of the problemwise effects.

\begin{table*}[t]
\centering
\caption{IGD comparison on 16 10-objective problems with 300 evaluations.}
\label{tab:exp:main10}
\scriptsize
\setlength{\tabcolsep}{3pt}
\renewcommand{\arraystretch}{1.12}
\begin{tabular}{lcccccccc}
\toprule
Problem & REMO & PIEA & MCEAD & R2AEA & PCSAEA & KRVEA & CSEA & HPDC-MaOEA \\
\midrule
DTLZ1 & \shortstack{$\mathbf{2.0748e+2}^{=}$\\$(4.35e+1)$} & \shortstack{$2.4207e+2^{=}$\\$(7.32e+1)$} & \shortstack{$2.3531e+2^{=}$\\$(4.54e+1)$} & \shortstack{$2.7206e+2^{-}$\\$(2.49e+1)$} & \shortstack{$3.6855e+2^{-\dagger}$\\$(2.80e+1)$} & \shortstack{$3.6855e+2^{-\dagger}$\\$(2.80e+1)$} & \shortstack{$2.2281e+2^{=}$\\$(4.96e+1)$} & \shortstack{$2.2944e+2$\\$(3.62e+1)$} \\
DTLZ2 & \shortstack{$1.0902e+0^{=}$\\$(9.26e-2)$} & \shortstack{$\mathbf{9.6871e-1}^{+}$\\$(5.69e-2)$} & \shortstack{$1.0974e+0^{=}$\\$(6.08e-2)$} & \shortstack{$1.4611e+0^{-}$\\$(5.68e-2)$} & \shortstack{$1.5378e+0^{-}$\\$(6.87e-2)$} & \shortstack{$1.5416e+0^{-}$\\$(6.16e-2)$} & \shortstack{$1.3229e+0^{-}$\\$(9.81e-2)$} & \shortstack{$1.0946e+0$\\$(8.88e-2)$} \\
DTLZ3 & \shortstack{$8.0813e+2^{+}$\\$(1.50e+2)$} & \shortstack{$\mathbf{5.5459e+2}^{+}$\\$(1.06e+2)$} & \shortstack{$5.7499e+2^{+}$\\$(1.31e+2)$} & \shortstack{$5.6353e+2^{+}$\\$(6.65e+1)$} & \shortstack{$1.3393e+3^{-}$\\$(8.67e+1)$} & \shortstack{$1.3574e+3^{-}$\\$(1.10e+2)$} & \shortstack{$9.0082e+2^{=}$\\$(1.72e+2)$} & \shortstack{$9.3919e+2$\\$(1.40e+2)$} \\
DTLZ4 & \shortstack{$\mathbf{1.1487e+0}^{=}$\\$(1.05e-1)$} & \shortstack{$1.2452e+0^{=}$\\$(4.05e-2)$} & \shortstack{$1.2487e+0^{-}$\\$(1.91e-2)$} & \shortstack{$1.5944e+0^{-}$\\$(7.20e-2)$} & \shortstack{$1.7339e+0^{-}$\\$(5.89e-2)$} & \shortstack{$1.7361e+0^{-}$\\$(6.76e-2)$} & \shortstack{$1.2198e+0^{=}$\\$(8.16e-2)$} & \shortstack{$1.1814e+0$\\$(1.12e-1)$} \\
DTLZ5 & \shortstack{$6.2052e-1^{=}$\\$(1.47e-1)$} & \shortstack{$\mathbf{2.9016e-1}^{+}$\\$(3.54e-2)$} & \shortstack{$3.7432e-1^{+}$\\$(6.73e-2)$} & \shortstack{$9.8650e-1^{-}$\\$(8.39e-2)$} & \shortstack{$9.8455e-1^{-}$\\$(6.03e-2)$} & \shortstack{$9.9808e-1^{-}$\\$(4.84e-2)$} & \shortstack{$7.6415e-1^{-}$\\$(1.07e-1)$} & \shortstack{$6.6049e-1$\\$(1.47e-1)$} \\
DTLZ6 & \shortstack{$1.6378e+1^{-}$\\$(7.30e-1)$} & \shortstack{$\mathbf{5.4851e+0}^{+}$\\$(1.81e+0)$} & \shortstack{$8.9744e+0^{+}$\\$(1.77e+0)$} & \shortstack{$1.3116e+1^{+}$\\$(1.22e+0)$} & \shortstack{$1.8023e+1^{-}$\\$(1.70e-1)$} & \shortstack{$1.8042e+1^{-}$\\$(1.58e-1)$} & \shortstack{$1.6064e+1^{-}$\\$(6.81e-1)$} & \shortstack{$1.4463e+1$\\$(1.10e+0)$} \\
DTLZ7 & \shortstack{$2.1252e+1^{-}$\\$(2.44e+0)$} & \shortstack{$1.8588e+1^{-}$\\$(5.79e+0)$} & \shortstack{$2.0445e+1^{-}$\\$(4.67e+0)$} & \shortstack{$\mathbf{6.2038e+0}^{+}$\\$(1.22e+0)$} & \shortstack{$2.9100e+1^{-}$\\$(1.69e+0)$} & \shortstack{$2.8736e+1^{-}$\\$(1.88e+0)$} & \shortstack{$2.7611e+1^{-}$\\$(2.24e+0)$} & \shortstack{$1.2122e+1$\\$(2.80e+0)$} \\
WFG1 & \shortstack{$3.2354e+0^{=}$\\$(8.74e-2)$} & \shortstack{$3.3433e+0^{-}$\\$(4.96e-2)$} & \shortstack{$3.4587e+0^{-}$\\$(3.65e-2)$} & \shortstack{$3.1986e+0^{=}$\\$(3.20e-2)$} & \shortstack{$3.4196e+0^{-}$\\$(1.68e-2)$} & \shortstack{$3.4256e+0^{-}$\\$(2.10e-2)$} & \shortstack{$3.2723e+0^{-}$\\$(6.63e-2)$} & \shortstack{$\mathbf{3.1933e+0}$\\$(2.79e-2)$} \\
WFG2 & \shortstack{$3.5108e+0^{-}$\\$(1.08e+0)$} & \shortstack{$2.3710e+0^{-}$\\$(4.69e-1)$} & \shortstack{$3.4836e+0^{-}$\\$(9.77e-1)$} & \shortstack{$2.2469e+0^{-}$\\$(1.13e-1)$} & \shortstack{$3.5949e+0^{-}$\\$(6.30e-1)$} & \shortstack{$3.8780e+0^{-}$\\$(6.00e-1)$} & \shortstack{$3.5554e+0^{-}$\\$(1.03e+0)$} & \shortstack{$\mathbf{2.1211e+0}$\\$(6.82e-1)$} \\
WFG3 & \shortstack{$1.5188e+0^{+}$\\$(1.01e-1)$} & \shortstack{$\mathbf{1.1823e+0}^{+}$\\$(1.49e-1)$} & \shortstack{$1.2551e+0^{+}$\\$(4.56e-2)$} & \shortstack{$1.5946e+0^{+}$\\$(4.99e-2)$} & \shortstack{$1.6472e+0^{+}$\\$(3.93e-2)$} & \shortstack{$1.6496e+0^{+}$\\$(4.64e-2)$} & \shortstack{$1.5723e+0^{+}$\\$(8.05e-2)$} & \shortstack{$1.7006e+0$\\$(5.87e-2)$} \\
WFG4 & \shortstack{$8.1770e+0^{-}$\\$(9.93e-1)$} & \shortstack{$\mathbf{5.7766e+0}^{+}$\\$(5.86e-1)$} & \shortstack{$7.1626e+0^{=}$\\$(8.70e-1)$} & \shortstack{$5.7901e+0^{+}$\\$(3.16e-1)$} & \shortstack{$9.1390e+0^{-}$\\$(4.55e-1)$} & \shortstack{$9.2700e+0^{-}$\\$(3.98e-1)$} & \shortstack{$9.0803e+0^{-}$\\$(8.62e-1)$} & \shortstack{$6.6311e+0$\\$(6.38e-1)$} \\
WFG5 & \shortstack{$6.9444e+0^{-}$\\$(5.14e-1)$} & \shortstack{$\mathbf{5.7544e+0}^{=}$\\$(2.59e-1)$} & \shortstack{$6.6193e+0^{-}$\\$(3.95e-1)$} & \shortstack{$6.5834e+0^{-}$\\$(4.62e-1)$} & \shortstack{$7.0750e+0^{-}$\\$(1.97e-1)$} & \shortstack{$6.9475e+0^{-}$\\$(2.47e-1)$} & \shortstack{$7.1022e+0^{-}$\\$(4.09e-1)$} & \shortstack{$5.9732e+0$\\$(4.00e-1)$} \\
WFG6 & \shortstack{$6.4255e+0^{-}$\\$(4.79e-1)$} & \shortstack{$5.9360e+0^{-}$\\$(4.03e-1)$} & \shortstack{$7.2512e+0^{-}$\\$(6.73e-1)$} & \shortstack{$\mathbf{5.2560e+0}^{=}$\\$(1.32e-1)$} & \shortstack{$7.3710e+0^{-}$\\$(3.74e-1)$} & \shortstack{$7.3224e+0^{-}$\\$(3.02e-1)$} & \shortstack{$7.4553e+0^{-}$\\$(5.14e-1)$} & \shortstack{$5.2773e+0$\\$(2.01e-1)$} \\
WFG7 & \shortstack{$6.2812e+0^{-}$\\$(7.42e-1)$} & \shortstack{$6.9139e+0^{-}$\\$(5.04e-1)$} & \shortstack{$8.1581e+0^{-}$\\$(6.29e-1)$} & \shortstack{$6.4875e+0^{-}$\\$(4.09e-1)$} & \shortstack{$7.8445e+0^{-}$\\$(3.81e-1)$} & \shortstack{$7.8642e+0^{-}$\\$(4.20e-1)$} & \shortstack{$7.3904e+0^{-}$\\$(6.15e-1)$} & \shortstack{$\mathbf{5.2969e+0}$\\$(2.57e-1)$} \\
WFG8 & \shortstack{$7.0559e+0^{-}$\\$(5.06e-1)$} & \shortstack{$6.3573e+0^{-}$\\$(4.99e-1)$} & \shortstack{$7.4671e+0^{-}$\\$(4.72e-1)$} & \shortstack{$\mathbf{5.5601e+0}^{=}$\\$(1.19e-1)$} & \shortstack{$7.5364e+0^{-}$\\$(2.73e-1)$} & \shortstack{$7.4984e+0^{-}$\\$(3.01e-1)$} & \shortstack{$7.7877e+0^{-}$\\$(3.46e-1)$} & \shortstack{$5.6930e+0$\\$(2.66e-1)$} \\
WFG9 & \shortstack{$7.6476e+0^{-}$\\$(7.61e-1)$} & \shortstack{$\mathbf{5.7930e+0}^{+}$\\$(3.72e-1)$} & \shortstack{$7.0825e+0^{-}$\\$(6.96e-1)$} & \shortstack{$6.4412e+0^{=}$\\$(3.15e-1)$} & \shortstack{$7.7853e+0^{-}$\\$(4.44e-1)$} & \shortstack{$7.7844e+0^{-}$\\$(3.65e-1)$} & \shortstack{$8.1835e+0^{-}$\\$(3.07e-1)$} & \shortstack{$6.3788e+0$\\$(6.11e-1)$} \\
\midrule
W/T/L & 9/5/2 & 6/3/7 & 9/3/4 & 7/4/5 & 15/0/1 & 15/0/1 & 12/3/1 & -- \\
\bottomrule
\end{tabular}
\par\smallskip
\begin{minipage}{\linewidth}
\scriptsize Each cell contains the exported mean IGD and standard deviation on separate lines. Symbols describe the comparator relative to HPDC-MaOEA; W/T/L describes HPDC-MaOEA relative to the comparator. Bold denotes the smallest mean. $\dagger$: the two identical DTLZ1 entries are preserved from the source and require provenance verification.
\end{minipage}
\end{table*}

\begin{table*}[t]
\centering
\caption{IGD comparison on 16 20-objective problems with 300 evaluations.}
\label{tab:exp:main20}
\scriptsize
\setlength{\tabcolsep}{3pt}
\renewcommand{\arraystretch}{1.12}
\begin{tabular}{lcccccccc}
\toprule
Problem & REMO & PIEA & MCEAD & R2AEA & KRVEA & CSEA & HES-EA & HPDC-MaOEA \\
\midrule
DTLZ1 & \shortstack{$\mathbf{6.7408e+1}^{+}$\\$(1.93e+1)$} & \shortstack{$9.3984e+1^{-}$\\$(1.82e+1)$} & \shortstack{$7.6364e+1^{=}$\\$(3.61e+1)$} & \shortstack{$1.0965e+2^{-}$\\$(1.62e+1)$} & \shortstack{$1.5756e+2^{-}$\\$(1.98e+1)$} & \shortstack{$7.9813e+1^{=}$\\$(1.99e+1)$} & \shortstack{$1.6595e+2^{-}$\\$(1.85e+1)$} & \shortstack{$7.9755e+1$\\$(1.09e+1)$} \\
DTLZ2 & \shortstack{$\mathbf{1.1173e+0}^{=}$\\$(5.45e-2)$} & \shortstack{$1.1636e+0^{-}$\\$(2.17e-2)$} & \shortstack{$1.2196e+0^{-}$\\$(2.96e-2)$} & \shortstack{$1.3157e+0^{-}$\\$(4.47e-2)$} & \shortstack{$1.3465e+0^{-}$\\$(2.76e-2)$} & \shortstack{$1.1817e+0^{-}$\\$(5.61e-2)$} & \shortstack{$1.3478e+0^{-}$\\$(3.70e-2)$} & \shortstack{$1.1350e+0$\\$(4.89e-2)$} \\
DTLZ3 & \shortstack{$2.6646e+2^{=}$\\$(7.13e+1)$} & \shortstack{$2.5997e+2^{=}$\\$(1.15e+1)$} & \shortstack{$\mathbf{1.2168e+2}^{+}$\\$(1.12e+2)$} & \shortstack{$3.0721e+2^{=}$\\$(5.93e+1)$} & \shortstack{$5.1789e+2^{-}$\\$(7.55e+1)$} & \shortstack{$2.5514e+2^{=}$\\$(6.82e+1)$} & \shortstack{$5.4369e+2^{-}$\\$(8.04e+1)$} & \shortstack{$2.8568e+2$\\$(6.84e+1)$} \\
DTLZ4 & \shortstack{$\mathbf{1.0368e+0}^{=}$\\$(4.26e-2)$} & \shortstack{$1.1269e+0^{-}$\\$(6.03e-2)$} & \shortstack{$1.2852e+0^{-}$\\$(3.78e-2)$} & \shortstack{$1.2275e+0^{-}$\\$(5.62e-2)$} & \shortstack{$1.2558e+0^{-}$\\$(3.34e-2)$} & \shortstack{$1.0610e+0^{=}$\\$(4.81e-2)$} & \shortstack{$1.2647e+0^{-}$\\$(2.73e-2)$} & \shortstack{$1.0578e+0$\\$(5.34e-2)$} \\
DTLZ5 & \shortstack{$2.7068e-1^{=}$\\$(5.25e-2)$} & \shortstack{$1.5237e-1^{+}$\\$(2.26e-2)$} & \shortstack{$\mathbf{1.5161e-1}^{+}$\\$(4.26e-2)$} & \shortstack{$4.4572e-1^{-}$\\$(7.26e-2)$} & \shortstack{$4.1362e-1^{-}$\\$(4.70e-2)$} & \shortstack{$2.9000e-1^{=}$\\$(5.36e-2)$} & \shortstack{$4.1158e-1^{-}$\\$(3.61e-2)$} & \shortstack{$3.0345e-1$\\$(6.65e-2)$} \\
DTLZ6 & \shortstack{$7.3417e+0^{=}$\\$(6.33e-1)$} & \shortstack{$\mathbf{2.1640e+0}^{+}$\\$(1.05e+0)$} & \shortstack{$4.1169e+0^{+}$\\$(1.03e+0)$} & \shortstack{$5.6087e+0^{=}$\\$(1.98e+0)$} & \shortstack{$9.2234e+0^{-}$\\$(1.53e-1)$} & \shortstack{$7.4321e+0^{-}$\\$(7.00e-1)$} & \shortstack{$9.1974e+0^{-}$\\$(1.22e-1)$} & \shortstack{$6.6546e+0$\\$(9.39e-1)$} \\
DTLZ7 & \shortstack{$3.8239e+1^{-}$\\$(7.37e+0)$} & \shortstack{$1.8728e+1^{-}$\\$(9.86e+0)$} & \shortstack{$2.8155e+1^{-}$\\$(1.25e+1)$} & \shortstack{$\mathbf{1.4039e+1}^{+}$\\$(2.89e+0)$} & \shortstack{$4.9744e+1^{-}$\\$(5.36e+0)$} & \shortstack{$4.8818e+1^{-}$\\$(9.05e+0)$} & \shortstack{$4.8867e+1^{-}$\\$(4.77e+0)$} & \shortstack{$1.8154e+1$\\$(2.61e+0)$} \\
WFG1 & \shortstack{$5.6202e+0^{=}$\\$(8.82e-2)$} & \shortstack{$5.7081e+0^{-}$\\$(5.05e-2)$} & \shortstack{$5.7957e+0^{-}$\\$(4.51e-2)$} & \shortstack{$\mathbf{5.6059e+0}^{=}$\\$(4.65e-2)$} & \shortstack{$5.7218e+0^{-}$\\$(3.08e-2)$} & \shortstack{$5.6736e+0^{-}$\\$(1.06e-1)$} & \shortstack{$5.7151e+0^{-}$\\$(3.91e-2)$} & \shortstack{$5.6075e+0$\\$(7.93e-2)$} \\
WFG2 & \shortstack{$7.9480e+0^{-}$\\$(1.54e+0)$} & \shortstack{$6.0506e+0^{-}$\\$(1.05e+0)$} & \shortstack{$7.6391e+0^{-}$\\$(2.05e+0)$} & \shortstack{$\mathbf{4.6261e+0}^{+}$\\$(1.67e-1)$} & \shortstack{$8.3973e+0^{-}$\\$(9.72e-1)$} & \shortstack{$8.8363e+0^{-}$\\$(1.30e+0)$} & \shortstack{$8.8430e+0^{-}$\\$(1.37e+0)$} & \shortstack{$5.4393e+0$\\$(1.31e+0)$} \\
WFG3 & \shortstack{$2.2519e+0^{+}$\\$(2.32e-1)$} & \shortstack{$\mathbf{1.7073e+0}^{+}$\\$(1.73e-1)$} & \shortstack{$1.9974e+0^{+}$\\$(2.03e-1)$} & \shortstack{$2.6029e+0^{+}$\\$(1.15e-1)$} & \shortstack{$2.4955e+0^{+}$\\$(8.00e-2)$} & \shortstack{$2.3113e+0^{+}$\\$(1.84e-1)$} & \shortstack{$2.4573e+0^{+}$\\$(6.25e-2)$} & \shortstack{$2.7306e+0$\\$(1.52e-1)$} \\
WFG4 & \shortstack{$2.9066e+1^{-}$\\$(1.21e+0)$} & \shortstack{$2.0639e+1^{+}$\\$(3.25e+0)$} & \shortstack{$2.0456e+1^{+}$\\$(3.65e+0)$} & \shortstack{$\mathbf{1.7947e+1}^{+}$\\$(1.25e+0)$} & \shortstack{$2.9490e+1^{-}$\\$(7.24e-1)$} & \shortstack{$2.9600e+1^{-}$\\$(1.48e+0)$} & \shortstack{$2.9432e+1^{-}$\\$(9.66e-1)$} & \shortstack{$2.3927e+1$\\$(1.07e+0)$} \\
WFG5 & \shortstack{$2.5608e+1^{-}$\\$(1.19e+0)$} & \shortstack{$\mathbf{2.1972e+1}^{+}$\\$(9.12e-1)$} & \shortstack{$2.2931e+1^{=}$\\$(6.52e-1)$} & \shortstack{$2.3519e+1^{=}$\\$(1.28e+0)$} & \shortstack{$2.4741e+1^{-}$\\$(6.96e-1)$} & \shortstack{$2.5912e+1^{-}$\\$(9.45e-1)$} & \shortstack{$2.4824e+1^{-}$\\$(5.85e-1)$} & \shortstack{$2.2849e+1$\\$(1.62e+0)$} \\
WFG6 & \shortstack{$2.4041e+1^{-}$\\$(1.66e+0)$} & \shortstack{$2.2696e+1^{-}$\\$(1.82e+0)$} & \shortstack{$2.5327e+1^{-}$\\$(1.94e+0)$} & \shortstack{$\mathbf{1.5679e+1}^{+}$\\$(5.02e-1)$} & \shortstack{$2.5293e+1^{-}$\\$(5.92e-1)$} & \shortstack{$2.6268e+1^{-}$\\$(1.20e+0)$} & \shortstack{$2.5393e+1^{-}$\\$(7.75e-1)$} & \shortstack{$1.9097e+1$\\$(6.18e-1)$} \\
WFG7 & \shortstack{$2.4322e+1^{-}$\\$(1.10e+0)$} & \shortstack{$2.4579e+1^{-}$\\$(8.74e-1)$} & \shortstack{$2.7355e+1^{-}$\\$(1.36e+0)$} & \shortstack{$2.0597e+1^{=}$\\$(2.63e+0)$} & \shortstack{$2.6465e+1^{-}$\\$(7.60e-1)$} & \shortstack{$2.6430e+1^{-}$\\$(1.40e+0)$} & \shortstack{$2.5957e+1^{-}$\\$(7.48e-1)$} & \shortstack{$\mathbf{2.0592e+1}$\\$(7.61e-1)$} \\
WFG8 & \shortstack{$2.4953e+1^{-}$\\$(1.02e+0)$} & \shortstack{$2.3237e+1^{-}$\\$(1.24e+0)$} & \shortstack{$2.6026e+1^{-}$\\$(1.81e+0)$} & \shortstack{$\mathbf{1.6282e+1}^{+}$\\$(3.72e-1)$} & \shortstack{$2.5519e+1^{-}$\\$(4.89e-1)$} & \shortstack{$2.7053e+1^{-}$\\$(7.42e-1)$} & \shortstack{$2.5791e+1^{-}$\\$(6.48e-1)$} & \shortstack{$2.0738e+1$\\$(8.26e-1)$} \\
WFG9 & \shortstack{$2.6155e+1^{-}$\\$(1.38e+0)$} & \shortstack{$2.2604e+1^{=}$\\$(1.32e+0)$} & \shortstack{$2.3089e+1^{=}$\\$(1.39e+0)$} & \shortstack{$\mathbf{2.2569e+1}^{=}$\\$(1.27e+0)$} & \shortstack{$2.5820e+1^{-}$\\$(7.55e-1)$} & \shortstack{$2.7847e+1^{-}$\\$(8.01e-1)$} & \shortstack{$2.6166e+1^{-}$\\$(8.22e-1)$} & \shortstack{$2.3053e+1$\\$(1.30e+0)$} \\
\midrule
W/T/L & 8/6/2 & 9/2/5 & 8/3/5 & 4/6/6 & 15/0/1 & 11/4/1 & 15/0/1 & -- \\
\bottomrule
\end{tabular}
\par\smallskip
\begin{minipage}{\linewidth}
\scriptsize Each cell contains the exported mean IGD and standard deviation on separate lines. Symbols describe the comparator relative to HPDC-MaOEA; W/T/L describes HPDC-MaOEA relative to the comparator. Bold denotes the smallest mean.
\end{minipage}
\end{table*}

\paragraph{Additional objective counts.}
Table~\ref{tab:exp:objectives} summarises the common-baseline comparisons.
On the full 16-problem suite, the reported win fraction increased from
34.4\% at $M=3$ to 62.5\% at $M=20$. This is a descriptive comparison
across the recorded settings: the decision dimension also changes between
the lower- and higher-objective experiments. The $M=12$ result covered
only nine WFG problems and is therefore shown separately. It yielded
27/3/6 against the common baselines and 43/3/8 against all six available
baselines.

\begin{table}[t]
\centering
\caption{Common-baseline comparisons across objective counts.}
\label{tab:exp:objectives}
\scriptsize
\setlength{\tabcolsep}{2.5pt}
\renewcommand{\arraystretch}{1.12}
\begin{tabular}{lcccc}
\toprule
$M$ & Problems & W/T/L & Win fraction & Median reduction \\
\midrule
3 & 16 & 22/18/24 & 34.4\% & 2.06\% \\
5 & 16 & 31/18/15 & 48.4\% & 6.76\% \\
8 & 16 & 38/12/14 & 59.4\% & 7.04\% \\
10 & 16 & 39/11/14 & 60.9\% & 10.49\% \\
20 & 16 & 40/11/13 & 62.5\% & 9.73\% \\
\midrule
12 & 9 & 27/3/6 & 75.0\% & 15.66\% \\
\bottomrule
\end{tabular}
\par\smallskip
\begin{minipage}{\linewidth}
\scriptsize REMO, PIEA, MCEAD and KRVEA are included in every row. The $M=12$ row covers WFG1--WFG9 only. Settings and verification notes follow Tables~\ref{tab:exp:settings}--\ref{tab:exp:main20}.
\end{minipage}
\end{table}

\subsection{Ablation Study}
\label{sec:exp:ablation}

The module ablation compared REMO, REMO with hybrid PBI grouping (HPC),
REMO with the candidate module, and the full method. Because the number
of representatives also affects the host algorithm, the available
configurations included REMO and REMO+HPC at both $k=6$ and the paper
setting $k=15$ for $M=10$ or $k=30$ for $M=20$. The candidate-only
configuration used $k=6$. Tables~\ref{tab:exp:abl10}
and~\ref{tab:exp:abl20} give all six identified configurations. The
20-objective table contains 15 valid problems; its WFG4 row is left blank
because the source ablation export labels that instance $M=10$.

At a matched representative count, adding candidate selection to
REMO+HPC produced reported W/T/L counts of 3/12/1 at $M=10$ and 3/11/1
at $M=20$. On DTLZ7, mean IGD decreased from 17.305 to 12.122 and from
25.637 to 18.154, respectively. WFG9 improved at both objective counts;
the remaining reported wins were WFG7 at $M=10$ and WFG6 at $M=20$.
Thus, the matched comparison identifies gains on specific problems while
most differences were not detected by the exported pairwise tests.

The full method attained the lowest mean IGD in 7 of the 16
10-objective instances and 5 of the 15 valid 20-objective instances among
the six configurations. Its comparison with enlarged-$k$ REMO yielded
10/5/1 and 10/4/1. These latter contrasts measure the combined change
from both modules at a fixed representative count. In contrast, the
candidate-only $k=6$ configuration differs from the full method in both
grouping and $k$, and does not isolate grouping removal at the paper
setting.

\begin{table*}[t]
\centering
\caption{Module ablation at $M=10$ and 300 evaluations.}
\label{tab:exp:abl10}
\scriptsize
\setlength{\tabcolsep}{4pt}
\renewcommand{\arraystretch}{1.12}
\begin{tabular}{lcccccc}
\toprule
Problem & \shortstack{REMO\\$k=6$} & \shortstack{REMO\\$k=15$} & \shortstack{REMO+HPC\\$k=6$} & \shortstack{REMO+HPC\\$k=15$} & \shortstack{REMO+candidate\\$k=6$} & \shortstack{HPDC-MaOEA\\$k=15$} \\
\midrule
DTLZ1 & \shortstack{$2.0748e+2^{=}$\\$(4.35e+1)$} & \shortstack{$2.5572e+2^{-}$\\$(3.72e+1)$} & \shortstack{$\mathbf{1.9210e+2}^{+}$\\$(4.29e+1)$} & \shortstack{$2.2265e+2^{=}$\\$(2.72e+1)$} & \shortstack{$2.1327e+2^{=}$\\$(4.33e+1)$} & \shortstack{$2.2944e+2$\\$(3.62e+1)$} \\
DTLZ2 & \shortstack{$1.0902e+0^{=}$\\$(9.26e-2)$} & \shortstack{$1.1321e+0^{=}$\\$(8.14e-2)$} & \shortstack{$1.1179e+0^{=}$\\$(8.60e-2)$} & \shortstack{$1.1325e+0^{=}$\\$(1.03e-1)$} & \shortstack{$\mathbf{1.0434e+0}^{=}$\\$(1.08e-1)$} & \shortstack{$1.0946e+0$\\$(8.88e-2)$} \\
DTLZ3 & \shortstack{$8.0813e+2^{+}$\\$(1.50e+2)$} & \shortstack{$9.8158e+2^{=}$\\$(1.46e+2)$} & \shortstack{$\mathbf{7.4833e+2}^{+}$\\$(1.57e+2)$} & \shortstack{$9.4062e+2^{=}$\\$(1.36e+2)$} & \shortstack{$8.3489e+2^{+}$\\$(1.38e+2)$} & \shortstack{$9.3919e+2$\\$(1.40e+2)$} \\
DTLZ4 & \shortstack{$\mathbf{1.1487e+0}^{=}$\\$(1.05e-1)$} & \shortstack{$1.2536e+0^{=}$\\$(1.07e-1)$} & \shortstack{$1.1609e+0^{=}$\\$(1.08e-1)$} & \shortstack{$1.1720e+0^{=}$\\$(8.09e-2)$} & \shortstack{$1.1829e+0^{=}$\\$(1.08e-1)$} & \shortstack{$1.1814e+0$\\$(1.12e-1)$} \\
DTLZ5 & \shortstack{$6.2052e-1^{=}$\\$(1.47e-1)$} & \shortstack{$6.5860e-1^{=}$\\$(1.07e-1)$} & \shortstack{$\mathbf{5.2409e-1}^{+}$\\$(1.16e-1)$} & \shortstack{$6.0132e-1^{=}$\\$(1.35e-1)$} & \shortstack{$5.8917e-1^{=}$\\$(1.12e-1)$} & \shortstack{$6.6049e-1$\\$(1.47e-1)$} \\
DTLZ6 & \shortstack{$1.6378e+1^{-}$\\$(7.30e-1)$} & \shortstack{$1.6101e+1^{-}$\\$(6.03e-1)$} & \shortstack{$1.5155e+1^{=}$\\$(6.93e-1)$} & \shortstack{$1.5000e+1^{=}$\\$(5.75e-1)$} & \shortstack{$1.6196e+1^{-}$\\$(5.75e-1)$} & \shortstack{$\mathbf{1.4463e+1}$\\$(1.10e+0)$} \\
DTLZ7 & \shortstack{$2.1252e+1^{-}$\\$(2.44e+0)$} & \shortstack{$2.1381e+1^{-}$\\$(2.48e+0)$} & \shortstack{$1.7953e+1^{-}$\\$(3.03e+0)$} & \shortstack{$1.7305e+1^{-}$\\$(1.70e+0)$} & \shortstack{$1.4917e+1^{-}$\\$(1.86e+0)$} & \shortstack{$\mathbf{1.2122e+1}$\\$(2.80e+0)$} \\
WFG1 & \shortstack{$3.2354e+0^{=}$\\$(8.74e-2)$} & \shortstack{$3.2346e+0^{-}$\\$(5.80e-2)$} & \shortstack{$3.1879e+0^{=}$\\$(5.68e-2)$} & \shortstack{$3.2000e+0^{=}$\\$(3.80e-2)$} & \shortstack{$\mathbf{3.1734e+0}^{=}$\\$(6.75e-2)$} & \shortstack{$3.1933e+0$\\$(2.79e-2)$} \\
WFG2 & \shortstack{$3.5108e+0^{-}$\\$(1.08e+0)$} & \shortstack{$2.7019e+0^{-}$\\$(8.69e-1)$} & \shortstack{$2.7960e+0^{-}$\\$(8.43e-1)$} & \shortstack{$2.1621e+0^{=}$\\$(6.46e-1)$} & \shortstack{$3.0646e+0^{-}$\\$(1.04e+0)$} & \shortstack{$\mathbf{2.1211e+0}$\\$(6.82e-1)$} \\
WFG3 & \shortstack{$\mathbf{1.5188e+0}^{+}$\\$(1.01e-1)$} & \shortstack{$1.5502e+0^{+}$\\$(8.35e-2)$} & \shortstack{$1.5760e+0^{+}$\\$(8.86e-2)$} & \shortstack{$1.6399e+0^{+}$\\$(5.18e-2)$} & \shortstack{$1.6296e+0^{+}$\\$(6.34e-2)$} & \shortstack{$1.7006e+0$\\$(5.87e-2)$} \\
WFG4 & \shortstack{$8.1770e+0^{-}$\\$(9.93e-1)$} & \shortstack{$7.2123e+0^{-}$\\$(7.13e-1)$} & \shortstack{$8.6573e+0^{-}$\\$(1.21e+0)$} & \shortstack{$7.1851e+0^{=}$\\$(9.13e-1)$} & \shortstack{$7.5869e+0^{-}$\\$(9.62e-1)$} & \shortstack{$\mathbf{6.6311e+0}$\\$(6.38e-1)$} \\
WFG5 & \shortstack{$6.9444e+0^{-}$\\$(5.14e-1)$} & \shortstack{$5.9358e+0^{=}$\\$(4.19e-1)$} & \shortstack{$6.6419e+0^{-}$\\$(6.91e-1)$} & \shortstack{$\mathbf{5.8962e+0}^{=}$\\$(3.37e-1)$} & \shortstack{$6.4368e+0^{-}$\\$(5.01e-1)$} & \shortstack{$5.9732e+0$\\$(4.00e-1)$} \\
WFG6 & \shortstack{$6.4255e+0^{-}$\\$(4.79e-1)$} & \shortstack{$5.4737e+0^{-}$\\$(2.95e-1)$} & \shortstack{$6.4196e+0^{-}$\\$(5.05e-1)$} & \shortstack{$5.3355e+0^{=}$\\$(3.10e-1)$} & \shortstack{$5.7033e+0^{-}$\\$(2.69e-1)$} & \shortstack{$\mathbf{5.2773e+0}$\\$(2.01e-1)$} \\
WFG7 & \shortstack{$6.2812e+0^{-}$\\$(7.42e-1)$} & \shortstack{$5.7119e+0^{-}$\\$(4.43e-1)$} & \shortstack{$6.6489e+0^{-}$\\$(6.28e-1)$} & \shortstack{$5.4193e+0^{-}$\\$(2.07e-1)$} & \shortstack{$5.7946e+0^{-}$\\$(3.58e-1)$} & \shortstack{$\mathbf{5.2969e+0}$\\$(2.57e-1)$} \\
WFG8 & \shortstack{$7.0559e+0^{-}$\\$(5.06e-1)$} & \shortstack{$6.0420e+0^{-}$\\$(3.51e-1)$} & \shortstack{$6.7588e+0^{-}$\\$(4.65e-1)$} & \shortstack{$\mathbf{5.6556e+0}^{=}$\\$(1.43e-1)$} & \shortstack{$6.3053e+0^{-}$\\$(3.64e-1)$} & \shortstack{$5.6930e+0$\\$(2.66e-1)$} \\
WFG9 & \shortstack{$7.6476e+0^{-}$\\$(7.61e-1)$} & \shortstack{$6.9496e+0^{-}$\\$(4.45e-1)$} & \shortstack{$7.6437e+0^{-}$\\$(6.71e-1)$} & \shortstack{$6.7894e+0^{-}$\\$(6.20e-1)$} & \shortstack{$7.1350e+0^{-}$\\$(5.95e-1)$} & \shortstack{$\mathbf{6.3788e+0}$\\$(6.11e-1)$} \\
\midrule
W/T/L & 9/5/2 & 10/5/1 & 8/4/4 & 3/12/1 & 9/5/2 & -- \\
Best mean & 2 & 0 & 3 & 2 & 2 & 7 \\
\bottomrule
\end{tabular}
\par\smallskip
\begin{minipage}{\linewidth}
\scriptsize Cells report mean IGD (standard deviation). Symbols and W/T/L follow the main tables. HPC denotes hybrid PBI grouping. The enlarged-$k$ HPC-only configuration isolates addition of candidate selection when compared with the full method.
\end{minipage}
\end{table*}

\begin{table*}[t]
\centering
\caption{Module ablation at $M=20$ and 300 evaluations.}
\label{tab:exp:abl20}
\scriptsize
\setlength{\tabcolsep}{4pt}
\renewcommand{\arraystretch}{1.12}
\begin{tabular}{lcccccc}
\toprule
Problem & \shortstack{REMO\\$k=6$} & \shortstack{REMO\\$k=30$} & \shortstack{REMO+HPC\\$k=6$} & \shortstack{REMO+HPC\\$k=30$} & \shortstack{REMO+candidate\\$k=6$} & \shortstack{HPDC-MaOEA\\$k=30$} \\
\midrule
DTLZ1 & \shortstack{$6.7408e+1^{+}$\\$(1.93e+1)$} & \shortstack{$1.0194e+2^{-}$\\$(2.14e+1)$} & \shortstack{$\mathbf{5.7687e+1}^{+}$\\$(1.80e+1)$} & \shortstack{$8.1313e+1^{=}$\\$(2.10e+1)$} & \shortstack{$6.0747e+1^{+}$\\$(1.62e+1)$} & \shortstack{$7.9755e+1$\\$(1.09e+1)$} \\
DTLZ2 & \shortstack{$1.1173e+0^{=}$\\$(5.45e-2)$} & \shortstack{$1.1484e+0^{=}$\\$(4.27e-2)$} & \shortstack{$1.1203e+0^{=}$\\$(5.61e-2)$} & \shortstack{$1.1979e+0^{-}$\\$(3.90e-2)$} & \shortstack{$\mathbf{1.0857e+0}^{+}$\\$(7.42e-2)$} & \shortstack{$1.1350e+0$\\$(4.89e-2)$} \\
DTLZ3 & \shortstack{$2.6646e+2^{=}$\\$(7.13e+1)$} & \shortstack{$3.7834e+2^{-}$\\$(8.48e+1)$} & \shortstack{$\mathbf{2.2547e+2}^{+}$\\$(5.49e+1)$} & \shortstack{$3.2741e+2^{=}$\\$(5.90e+1)$} & \shortstack{$2.5467e+2^{=}$\\$(7.79e+1)$} & \shortstack{$2.8568e+2$\\$(6.84e+1)$} \\
DTLZ4 & \shortstack{$1.0368e+0^{=}$\\$(4.26e-2)$} & \shortstack{$1.0608e+0^{=}$\\$(3.77e-2)$} & \shortstack{$1.0592e+0^{=}$\\$(4.75e-2)$} & \shortstack{$\mathbf{1.0199e+0}^{+}$\\$(4.38e-2)$} & \shortstack{$1.0634e+0^{=}$\\$(5.86e-2)$} & \shortstack{$1.0578e+0$\\$(5.34e-2)$} \\
DTLZ5 & \shortstack{$2.7068e-1^{=}$\\$(5.25e-2)$} & \shortstack{$3.3109e-1^{=}$\\$(4.78e-2)$} & \shortstack{$\mathbf{2.0792e-1}^{+}$\\$(5.04e-2)$} & \shortstack{$2.6667e-1^{=}$\\$(4.58e-2)$} & \shortstack{$2.8410e-1^{=}$\\$(4.78e-2)$} & \shortstack{$3.0345e-1$\\$(6.65e-2)$} \\
DTLZ6 & \shortstack{$7.3417e+0^{=}$\\$(6.33e-1)$} & \shortstack{$7.5855e+0^{-}$\\$(4.73e-1)$} & \shortstack{$\mathbf{6.3369e+0}^{=}$\\$(8.09e-1)$} & \shortstack{$6.5356e+0^{=}$\\$(7.58e-1)$} & \shortstack{$7.4829e+0^{-}$\\$(6.41e-1)$} & \shortstack{$6.6546e+0$\\$(9.39e-1)$} \\
DTLZ7 & \shortstack{$3.8239e+1^{-}$\\$(7.37e+0)$} & \shortstack{$4.2108e+1^{-}$\\$(7.08e+0)$} & \shortstack{$2.9216e+1^{-}$\\$(4.21e+0)$} & \shortstack{$2.5637e+1^{-}$\\$(5.00e+0)$} & \shortstack{$2.2735e+1^{-}$\\$(3.18e+0)$} & \shortstack{$\mathbf{1.8154e+1}$\\$(2.61e+0)$} \\
WFG1 & \shortstack{$5.6202e+0^{=}$\\$(8.82e-2)$} & \shortstack{$5.6481e+0^{-}$\\$(4.89e-2)$} & \shortstack{$5.5981e+0^{=}$\\$(1.05e-1)$} & \shortstack{$5.6149e+0^{=}$\\$(7.80e-2)$} & \shortstack{$\mathbf{5.5705e+0}^{=}$\\$(8.88e-2)$} & \shortstack{$5.6075e+0$\\$(7.93e-2)$} \\
WFG2 & \shortstack{$7.9480e+0^{-}$\\$(1.54e+0)$} & \shortstack{$7.0318e+0^{-}$\\$(1.58e+0)$} & \shortstack{$7.2512e+0^{-}$\\$(1.65e+0)$} & \shortstack{$\mathbf{5.1016e+0}^{=}$\\$(6.88e-1)$} & \shortstack{$6.8209e+0^{-}$\\$(2.06e+0)$} & \shortstack{$5.4393e+0$\\$(1.31e+0)$} \\
WFG3 & \shortstack{$\mathbf{2.2519e+0}^{+}$\\$(2.32e-1)$} & \shortstack{$2.3254e+0^{+}$\\$(1.85e-1)$} & \shortstack{$2.5155e+0^{+}$\\$(1.99e-1)$} & \shortstack{$2.6644e+0^{=}$\\$(1.17e-1)$} & \shortstack{$2.4822e+0^{+}$\\$(1.70e-1)$} & \shortstack{$2.7306e+0$\\$(1.52e-1)$} \\
WFG4 & -- & -- & -- & -- & -- & -- \\
WFG5 & \shortstack{$2.5608e+1^{-}$\\$(1.19e+0)$} & \shortstack{$2.3181e+1^{=}$\\$(1.06e+0)$} & \shortstack{$2.5693e+1^{-}$\\$(9.94e-1)$} & \shortstack{$2.3849e+1^{=}$\\$(1.38e+0)$} & \shortstack{$2.4357e+1^{-}$\\$(9.19e-1)$} & \shortstack{$\mathbf{2.2849e+1}$\\$(1.62e+0)$} \\
WFG6 & \shortstack{$2.4041e+1^{-}$\\$(1.66e+0)$} & \shortstack{$2.1493e+1^{-}$\\$(7.92e-1)$} & \shortstack{$2.3492e+1^{-}$\\$(1.36e+0)$} & \shortstack{$1.9109e+1^{=}$\\$(1.02e+0)$} & \shortstack{$2.0834e+1^{-}$\\$(9.11e-1)$} & \shortstack{$\mathbf{1.9097e+1}$\\$(6.18e-1)$} \\
WFG7 & \shortstack{$2.4322e+1^{-}$\\$(1.10e+0)$} & \shortstack{$2.2442e+1^{-}$\\$(1.43e+0)$} & \shortstack{$2.3916e+1^{-}$\\$(2.00e+0)$} & \shortstack{$\mathbf{2.0450e+1}^{=}$\\$(8.73e-1)$} & \shortstack{$2.2471e+1^{-}$\\$(2.16e+0)$} & \shortstack{$2.0592e+1$\\$(7.61e-1)$} \\
WFG8 & \shortstack{$2.4953e+1^{-}$\\$(1.02e+0)$} & \shortstack{$2.2718e+1^{-}$\\$(1.02e+0)$} & \shortstack{$2.4403e+1^{-}$\\$(1.29e+0)$} & \shortstack{$2.1098e+1^{=}$\\$(5.56e-1)$} & \shortstack{$2.2424e+1^{-}$\\$(1.08e+0)$} & \shortstack{$\mathbf{2.0738e+1}$\\$(8.26e-1)$} \\
WFG9 & \shortstack{$2.6155e+1^{-}$\\$(1.38e+0)$} & \shortstack{$2.4590e+1^{-}$\\$(1.41e+0)$} & \shortstack{$2.5690e+1^{-}$\\$(1.26e+0)$} & \shortstack{$2.4763e+1^{-}$\\$(1.35e+0)$} & \shortstack{$2.4951e+1^{-}$\\$(1.72e+0)$} & \shortstack{$\mathbf{2.3053e+1}$\\$(1.30e+0)$} \\
\midrule
W/T/L & 7/6/2 & 10/4/1 & 7/4/4 & 3/11/1 & 8/4/3 & -- \\
Best mean & 1 & 0 & 4 & 3 & 2 & 5 \\
\bottomrule
\end{tabular}
\par\smallskip
\begin{minipage}{\linewidth}
\scriptsize Cells report mean IGD (standard deviation). Symbols and W/T/L follow the main tables. HPC denotes hybrid PBI grouping. The enlarged-$k$ HPC-only configuration isolates addition of candidate selection when compared with the full method. WFG4 is unfilled because the source row records $M=10$; all counts use the remaining 15 problems.
\end{minipage}
\end{table*}

\TODO{Insert the candidate-only configuration at $k=15,30$ and complete
the matched component-removal comparisons. Leave unverified reduced-$k$
full-method entries blank until their implementation identity is resolved.
The joint interaction of the two modules remains to be tested.}

\subsection{Analysis of Hybrid PBI-Based Quality Grouping}
\label{sec:exp:hpc}

\paragraph{Positive-group composition.}
Hybrid PBI grouping selected a positive group distinct from either
single-view ranking. This was examined using the archived snapshot study
on DTLZ2, DTLZ4, DTLZ7, WFG3 and WFG7 at $M=10,20$, with five runs for
each of two search behaviours: Hybrid and the native representative-based
binary grouping. The study comprised 100 trajectories and 6,700
checkpoints. The same snapshots were used for offline comparisons of
positive-group definitions, each retaining the top 25\% except for the
native binary group.

Along the 50 Hybrid trajectories, mean Jaccard overlap was 0.5470 between
the Hybrid group and the direction-score group, and 0.3996 between the
Hybrid group and the representative-margin group. Along the 50
native-group trajectories, the corresponding values were 0.5699 and
0.4244. These run-averaged overlaps establish a change in group
composition under both recorded search behaviours. Here, the continuous
representative margin is an equal-quota comparator. The production
Hybrid score combines the direction score with the binary label, as
specified in Section~\ref{sec:method}; it does not combine two continuous
scores.

\paragraph{Association with subsequent retention.}
The completed good-group precision study evaluated the same five
problems at $M=10,20$ with 25 runs per configuration, giving 250
trajectories at $FE_{\max}=500$. At each checkpoint, the Hybrid,
direction-score and representative-margin rankings selected the top
25\% of the current population. Precision was the fraction of selected
solutions that remained in the final population of the recorded Hybrid
trajectory. Checkpoint precision was averaged within each run and one
of four $FE/FE_{\max}$ stages, followed by equal averaging across runs
and configurations.

\begin{table}[t]
\centering
\caption{Precision@25\% for final-population retention.}
\label{tab:exp:precision}
\footnotesize
\setlength{\tabcolsep}{3pt}
\renewcommand{\arraystretch}{1.12}
\begin{tabular}{lccc}
\toprule
$FE/FE_{\max}$ & Hybrid & Direction & Margin \\
\midrule
$[0,0.25]$ & 0.08547 & 0.05949 & 0.07949 \\
$(0.25,0.50]$ & 0.11804 & 0.10751 & 0.12414 \\
$(0.50,0.75]$ & 0.29011 & 0.28834 & 0.27126 \\
$(0.75,1.00]$ & 0.65431 & 0.65458 & 0.62980 \\
\midrule
All stages & 0.28698 & 0.27748 & 0.27617 \\
\bottomrule
\end{tabular}
\par\smallskip
\begin{minipage}{\linewidth}
\scriptsize Each stage averages 250 run-level values per view. The last row gives equal weight to the four stages. Larger values mean a greater fraction of selected solutions remained in the final population.
\end{minipage}
\end{table}

Hybrid achieved mean Precision@25\% of 0.28698, compared with 0.27748
for the direction score and 0.27617 for the representative margin
(Table~\ref{tab:exp:precision}). Its advantage was not uniform across
stages: the representative margin had the highest mean in the second
stage, while the direction score and Hybrid were nearly equal in the
last stage. The native binary group selected 49.05\% of the population
on average and had precision 0.31905; its different selection fraction
requires it to be reported separately from the top-25\% comparisons.

The two single-view groups also contained distinct future-retained
solutions. Across 40 problem--objective--stage cells, 33 passed the
Holm-adjusted test requiring both views to contribute unique
future-retained solutions in a majority of runs. The stricter test
requiring Hybrid precision to exceed both single-view precisions passed
in one cell. Thus, the evidence supports nonredundant information in the
two views and a modest aggregate retention advantage for Hybrid along
the recorded trajectories. The endpoint measured here is future
population membership, rather than final algorithm-level IGD.

\paragraph{Effect of the mixing schedule.}
The offline snapshot comparison additionally evaluated a fixed
equal-weight mixture and a reversed schedule. Along the Hybrid
trajectories, their mean positive-group overlaps with the original
schedule were 0.99987 and 0.98806, respectively. In these snapshots,
group membership was therefore largely preserved by the alternative
schedules. These observations complement the ordering rationale in the
method section; they do not establish an endpoint advantage for the
time-varying schedule.

\TODO{Add matched closed-loop comparisons of the grouping schedules and
independent utility results after the reference-stability check is resolved.}

\subsection{Analysis of Dual-Mode Candidate Selection}
\label{sec:exp:candidate}

\paragraph{Batch expenditure and selection criteria.}
The candidate-value probe compared five policies on DTLZ2, DTLZ7, WFG3
and WFG7 with $M=20$, $N=100$ and $FE_{\max}=300$. Each policy had ten
runs per problem, paired by the problem and run seed. V0 applied the
REMO selection rule within the probe host. V1 selected the six highest
relation scores from the accumulated pool. V2 used exploration-only
selection, V3 indicator-only selection, and V4 the full dual-mode
strategy. V0 is a transplanted selection rule, not the original REMO
algorithm. V1--V4 used the same accumulated-pool construction and a
six-candidate batch bound; their search trajectories generated their
own pools.

All 200 runs completed at 300 evaluations. V1--V4 selected exactly six
candidates in every non-truncated iteration, with the final batch
reduced to two by the remaining budget. Their mean batch size was
5.882. Under V0, the batch size ranged from 1 to 38 and averaged 10.386.
This comparison directly demonstrates the expenditure control supplied
by the explicit batch bound.

\paragraph{Batch properties and candidate utility.}
Batch spread was the mean pairwise Euclidean distance between selected
candidates after scaling each decision variable by its bounds. Retention was the fraction of evaluated candidates accepted
by the host's population-selection operation. The within-pool gain
ratio compared the selected batch's reduction in mean reference-to-set
distance with that of a greedy reference batch. This diagnostic used
300 objective-space reference points and a 400-candidate subset that
included all selected candidates. Its objective evaluations were used
only for offline assessment and were not returned to the search.
Retention and gain ratios below use iterations with
$FE/FE_{\max}\geq0.5$; batch spread and size use the full trajectory.

\begin{table*}[t]
\centering
\caption{Candidate-value probe at $M=20$: final IGD and batch properties.}
\label{tab:exp:candidate}
\scriptsize
\setlength{\tabcolsep}{4pt}
\renewcommand{\arraystretch}{1.12}
\begin{tabular}{lccccc}
\toprule
Problem / metric & V0: REMO rule & V1: relation top-six & V2: exploration & V3: indicator & V4: dual mode \\
\midrule
DTLZ2 & \shortstack{$1.2052e+00$\\$(6.00e-02)$} & \shortstack{$1.2185e+00$\\$(3.75e-02)$} & \shortstack{$1.2165e+00$\\$(4.32e-02)$} & \shortstack{$1.1124e+00$\\$(5.25e-02)$} & \shortstack{$\mathbf{1.1109e+00}$\\$(4.67e-02)$} \\
DTLZ7 & \shortstack{$2.6138e+01$\\$(4.68e+00)$} & \shortstack{$\mathbf{1.7364e+01}$\\$(5.53e+00)$} & \shortstack{$2.0498e+01$\\$(9.13e+00)$} & \shortstack{$3.2509e+01$\\$(6.80e+00)$} & \shortstack{$1.9675e+01$\\$(2.16e+00)$} \\
WFG3 & \shortstack{$2.7212e+00$\\$(1.18e-01)$} & \shortstack{$2.8720e+00$\\$(8.43e-02)$} & \shortstack{$2.8239e+00$\\$(1.10e-01)$} & \shortstack{$\mathbf{2.6946e+00}$\\$(1.63e-01)$} & \shortstack{$2.7877e+00$\\$(1.40e-01)$} \\
WFG7 & \shortstack{$2.0776e+01$\\$(6.51e-01)$} & \shortstack{$2.1052e+01$\\$(6.09e-01)$} & \shortstack{$2.1241e+01$\\$(1.04e+00)$} & \shortstack{$2.2036e+01$\\$(1.18e+00)$} & \shortstack{$\mathbf{2.0318e+01}$\\$(1.42e+00)$} \\
\midrule
Mean batch size & 10.386 & 5.882 & 5.882 & 5.882 & 5.882 \\
Mean batch spread & 1.806 & 1.637 & 2.031 & 1.031 & 1.479 \\
Late retention & 0.756 & 0.706 & 0.676 & 0.761 & 0.770 \\
Late gain ratio & 0.123 & 0.086 & 0.081 & 0.059 & 0.105 \\
\bottomrule
\end{tabular}
\par\smallskip
\begin{minipage}{\linewidth}
\scriptsize IGD entries are means (sample standard deviations) over ten runs per problem and policy. Batch statistics are equally averaged across the 40 runs of each policy. Only retention and gain ratio use the late-stage restriction. The V0 rule uses the probe host, and the greedy-reference diagnostic is restricted to the sampled pool.
\end{minipage}
\end{table*}

The two modes produced different batch properties. Exploration-only
selection had the largest mean spread (2.031), whereas indicator-only
selection had the smallest (1.031). V4 attained late-stage retention
of 0.770 and a gain ratio of 0.105, compared with 0.706 and 0.086 for
relation top-six selection. V4's mean spread was 1.479, between the two
single-mode values. These measurements describe the selected batches
and are interpreted jointly with the problemwise endpoint IGD in
Table~\ref{tab:exp:candidate}.

V4 had lower mean final IGD than V1 on DTLZ2, WFG3 and WFG7, but higher
mean IGD on DTLZ7 (19.675 versus 17.364). Indicator-only selection
illustrates why retention alone is insufficient: its late-stage
retention was 0.761, but its mean IGD on DTLZ7 was 32.509. The two
criteria therefore differ in both batch composition and final quality,
and the benefit of their combination depends on the problem.

\paragraph{Within-host routing comparison.}
A separate routing experiment compared exploration-only,
indicator-only, linearly scheduled switching and equal-probability
switching on all 16 problems at $M=10,20$. These results used the
\texttt{UniformMix} implementation of the routing-study host; they are
reported separately from the \texttt{UniformMix\_Original} main results.
Table~\ref{tab:exp:routing} retains the IGD and available IGD$^+$
comparisons within that host.

Equal-probability switching obtained reported IGD W/T/L counts of
5/9/2 against either single mode at $M=10$. At $M=20$, the counts were
4/11/1 against exploration-only and 5/11/0 against indicator-only.
The DTLZ7 mean IGD was 18.615 under equal-probability switching and
33.514 under indicator-only selection. Against linear switching, the
counts were 0/16/0 at $M=10$ and 1/15/0 at $M=20$. Thus, these results
support using both criteria within this host, while showing little
separation between the two switching schedules.

\begin{table}[t]
\centering
\caption{Equal-probability switching versus alternative routing policies in the routing-study host.}
\label{tab:exp:routing}
\scriptsize
\setlength{\tabcolsep}{3pt}
\renewcommand{\arraystretch}{1.12}
\begin{tabular}{lccc}
\toprule
$M$ & Indicator & Comparator & W/T/L \\
\midrule
10 & IGD & Exploration-only & 5/9/2 \\
10 & IGD & Indicator-only & 5/9/2 \\
10 & IGD & Linear switching & 0/16/0 \\
\midrule
20 & IGD & Exploration-only & 4/11/1 \\
20 & IGD & Indicator-only & 5/11/0 \\
20 & IGD & Linear switching & 1/15/0 \\
\midrule
20 & IGD$^+$ & Exploration-only & 3/12/1 \\
20 & IGD$^+$ & Indicator-only & 5/9/2 \\
20 & IGD$^+$ & Linear switching & 1/15/0 \\
\bottomrule
\end{tabular}
\par\smallskip
\begin{minipage}{\linewidth}
\scriptsize Counts use all 16 problems and the exported annotations relative to UniformMix. These are separate host-specific comparisons; they are not ablations of the UniformMix\_Original entries in the main tables.
\end{minipage}
\end{table}

\TODO{Complete the relation-top-six and single-mode ablations for the
final main-comparison implementation, together with the missing
10-objective candidate-value probe.}

\subsection{Robustness to Core Controls}
\label{sec:exp:sensitivity}

\paragraph{Stability of positive-group membership.}
The offline grouping study tested removal of 5\% and 10\% of the
population before reconstructing the group. Mean retained-set Jaccard
overlap for Hybrid was 0.97074 and 0.95471 along the Hybrid trajectories,
respectively. The native-group trajectories yielded 0.97081 and
0.95470. Group membership was therefore stable under these local
population perturbations. This test concerns label stability at fixed
algorithm controls.

\paragraph{Additional evaluation budget.}
The available 500-evaluation comparison covered the 16 10-objective
problems with the \texttt{UniformMix} host. Its reported W/T/L counts
were 10/4/2 against REMO, 7/4/5 against PIEA, 11/1/4 against MCEAD,
7/4/5 against R2AEA, 8/4/4 against HES-EA and 15/1/0 against PCSAEA.
This provides additional within-budget performance evidence. A
matched budget comparison for the final implementation remains open.

\TODO{Add final-implementation sensitivity results for the core
grouping and candidate-allocation controls, matched FE-budget comparisons,
and the wall-clock shares of the two modules.}

\clearpage
